% Nguồn: Tuyển tập đề tham khảo môn Toán TS10 2025-2026 – TP. Hồ Chí Minh
% Đợt 6: Quận 5 (Đề 3 - THTH Sài Gòn) và Quận 6 (Đề 1) (14 câu tự luận)

% ==============================================================================
% QUẬN 5 – ĐỀ 3 (TRƯỜNG TRUNG HỌC THỰC HÀNH SÀI GÒN)
% ==============================================================================

\begin{bt}%[9D6H18-2]
% nguon: Đề tham khảo Tuyển sinh 10 TP.HCM 2025-2026 – Quận 5 – Đề 3 – Bài 1
Cho parabol $(P)\colon y = \dfrac{x^2}{4}$.
\begin{enumerate}[a)]
    \item Vẽ đồ thị $(P)$ trên hệ trục tọa độ $Oxy$.
    \item Tìm tọa độ điểm $M$ thuộc $(P)$ có tung độ nhỏ hơn hoành độ 1 đơn vị.
\end{enumerate}
\loigiai{
\begin{enumerate}[a)]
    \item Bảng giá trị tương ứng của $x$ và $y$:
    \begin{center}
    \begin{tabular}{|c|c|c|c|c|c|}
        \hline
        $x$ & $-4$ & $-2$ & $0$ & $2$ & $4$ \\ \hline
        $y = \dfrac{x^2}{4}$ & $4$ & $1$ & $0$ & $1$ & $4$ \\ \hline
    \end{tabular}
    \end{center}
    Đồ thị $(P)$ là một parabol đi qua gốc tọa độ $O(0;0)$, nhận trục $Oy$ làm trục đối xứng và bề lõm quay lên phía trên.
    \item Điểm $M(x_M; y_M) \in (P)$ có tung độ nhỏ hơn hoành độ 1 đơn vị, ta có hệ thức:
    $y_M = x_M - 1$.\\
    Vì $M \in (P)$ nên ta có phương trình:
    $\dfrac{x_M^2}{4} = x_M - 1 \Leftrightarrow x_M^2 - 4x_M + 4 = 0 \Leftrightarrow (x_M - 2)^2 = 0 \Leftrightarrow x_M = 2$.\\
    Với $x_M = 2 \Rightarrow y_M = 2 - 1 = 1$.\\
    Vậy tọa độ điểm cần tìm là $M(2; 1)$.
\end{enumerate}
}
\end{bt}

\begin{bt}%[9D6H20-2]
% nguon: Đề tham khảo Tuyển sinh 10 TP.HCM 2025-2026 – Quận 5 – Đề 3 – Bài 2
Cho phương trình $4x^2 + 12x + 3 = 0$ \quad (1).
\begin{enumerate}[a)]
    \item Chứng minh phương trình (1) có hai nghiệm phân biệt.
    \item Gọi $x_1$, $x_2$ là hai nghiệm của phương trình (1). Không giải phương trình, tính giá trị của biểu thức $A = x_1^3 + x_2^3 + 2025$.
\end{enumerate}
\loigiai{
\begin{enumerate}[a)]
    \item Phương trình (1) có $a = 4$, $b' = 6$, $c = 3$.\\
    Biệt thức $\Delta' = (b')^2 - ac = 6^2 - 4 \cdot 3 = 36 - 12 = 24 > 0$.\\
    Vì $\Delta' > 0$ nên phương trình (1) luôn có hai nghiệm phân biệt $x_1, x_2$.
    \item Theo định lí Viète, ta có:
    $\left\{\begin{aligned} &S = x_1 + x_2 = -\dfrac{12}{4} = -3 \\ &P = x_1 x_2 = \dfrac{3}{4}. \end{aligned}\right.$\\
    Ta có hằng đẳng thức:
    $x_1^3 + x_2^3 = (x_1 + x_2)(x_1^2 - x_1 x_2 + x_2^2) = (x_1 + x_2)\left[(x_1 + x_2)^2 - 3x_1 x_2\right] = S(S^2 - 3P)$.\\
    Thay giá trị của $S$ và $P$ vào biểu thức $A$:
    $\begin{aligned}
    A &= S(S^2 - 3P) + 2025 = (-3)\left[(-3)^2 - 3 \cdot \dfrac{3}{4}\right] + 2025 \\
    &= (-3)\left(9 - \dfrac{9}{4}\right] + 2025 = (-3) \cdot \dfrac{27}{4} + 2025 = -\dfrac{81}{4} + 2025 = \dfrac{8019}{4} = 2004{,}75.
    \end{aligned}$\\
    Vậy $A = \dfrac{8019}{4} = 2004{,}75$.
\end{enumerate}
}
\end{bt}

\begin{bt}%[9D8H26-3]
% nguon: Đề tham khảo Tuyển sinh 10 TP.HCM 2025-2026 – Quận 5 – Đề 3 – Bài 3
Biểu đồ cột kép thống kê số học sinh khối 6 của trường THCS A tham gia các câu lạc bộ (mỗi bạn chỉ tham gia đúng một câu lạc bộ):
\begin{itemize}
    \item Cờ vua: 28 nam, 17 nữ.
    \item Dân ca: 15 nam, 45 nữ.
    \item Khiêu vũ: 24 nam, 40 nữ.
    \item Bơi lội: 35 nam, 24 nữ.
\end{itemize}
\begin{enumerate}[a)]
    \item Câu lạc bộ nào có sự chênh lệch nhiều nhất giữa số nam sinh và nữ sinh?
    \item Chọn ngẫu nhiên một học sinh khối 6, tính xác suất của các biến cố sau:
    \begin{itemize}
        \item $A$: “Học sinh được chọn là nữ”.
        \item $B$: “Học sinh được chọn không tham gia câu lạc bộ bơi lội và câu lạc bộ khiêu vũ”.
    \end{itemize}
\end{enumerate}
\loigiai{
\begin{enumerate}[a)]
    \item Chênh lệch giữa số học sinh nam và nữ tại từng câu lạc bộ là:
    \begin{itemize}
        \item Câu lạc bộ Cờ vua: $|28 - 17| = 11$ học sinh.
        \item Câu lạc bộ Dân ca: $|15 - 45| = 30$ học sinh.
        \item Câu lạc bộ Khiêu vũ: $|24 - 40| = 16$ học sinh.
        \item Câu lạc bộ Bơi lội: $|35 - 24| = 11$ học sinh.
    \end{itemize}
    Vậy câu lạc bộ Dân ca có sự chênh lệch nhiều nhất giữa số học sinh nam và nữ (chênh lệch 30 bạn).
    \item Tổng số học sinh khối 6 tham gia các câu lạc bộ là:
    $(28 + 17) + (15 + 45) + (24 + 40) + (35 + 24) = 45 + 60 + 64 + 59 = 228$ (học sinh).\\
    Do đó không gian mẫu có $n(\Omega) = 228$.
    \begin{itemize}
        \item Biến cố $A$: “Học sinh được chọn là nữ”.\\
        Tổng số học sinh nữ là: $17 + 45 + 40 + 24 = 126$ (học sinh).\\
        Số kết quả thuận lợi cho $A$ là $n(A) = 126$.\\
        Xác suất của biến cố $A$ là:
        $P(A) = \dfrac{126}{228} = \dfrac{21}{38}$.
        \item Biến cố $B$: “Học sinh được chọn không tham gia câu lạc bộ bơi lội và câu lạc bộ khiêu vũ”.\\
        Học sinh không tham gia bơi lội và khiêu vũ nghĩa là học sinh đó tham gia câu lạc bộ Cờ vua hoặc Dân ca.\\
        Số học sinh tham gia Cờ vua và Dân ca là:
        $45 + 60 = 105$ (học sinh).\\
        Số kết quả thuận lợi cho biến cố $B$ là $n(B) = 105$.\\
        Xác suất của biến cố $B$ là:
        $P(B) = \dfrac{105}{228} = \dfrac{35}{76}$.
    \end{itemize}
\end{enumerate}
}
\end{bt}

\begin{bt}%[9D6V21-3]
% nguon: Đề tham khảo Tuyển sinh 10 TP.HCM 2025-2026 – Quận 5 – Đề 3 – Bài 4
Chợ Dân Sinh chuẩn bị được đưa vào hoạt động. Nếu mỗi gian hàng của chợ này cho thuê với mức giá 31 triệu đồng/năm thì sẽ có 60 gian hàng được thuê. Theo khảo sát và nghiên cứu, nếu chợ giảm giá thuê và ước tính cứ mỗi lần giảm giá thuê 1 triệu đồng/năm thì số lượng gian hàng được thuê sẽ tăng thêm 20 gian hàng. Gọi $x$ (triệu đồng) là giá cho thuê mới ($0 < x < 31$).
\begin{enumerate}[a)]
    \item Viết biểu thức $P(x)$ biểu diễn theo $x$ doanh thu của chợ Dân Sinh.
    \item Phí duy trì mỗi gian hàng là 27 triệu đồng/năm. Hỏi giá cho thuê mới của chợ là bao nhiêu để lợi nhuận thu được là 245 triệu đồng? Biết lợi nhuận bằng doanh thu trừ đi tổng chi phí duy trì gian hàng.
\end{enumerate}
\loigiai{
\begin{enumerate}[a)]
    \item Khi giá cho thuê mới là $x$ (triệu đồng/năm), mức giảm giá so với giá ban đầu là:
    $31 - x$ (triệu đồng/năm).\\
    Cứ giảm 1 triệu đồng thì số gian hàng thuê tăng thêm 20 gian hàng, nên số gian hàng được thuê là:
    $N = 60 + 20(31 - x) = 60 + 620 - 20x = 680 - 20x$ (gian hàng).\\
    Doanh thu của chợ Dân Sinh theo giá thuê $x$ được biểu diễn bởi:
    $P(x) = x(680 - 20x) = -20x^2 + 680x$ (triệu đồng).
    \item Tổng chi phí duy trì của tất cả các gian hàng được thuê là:
    $C = 27(680 - 20x)$ (triệu đồng).\\
    Lợi nhuận của chợ được tính theo công thức:
    $\text{Lợi nhuận} = P(x) - C = x(680 - 20x) - 27(680 - 20x) = (x - 27)(680 - 20x)$.\\
    Để lợi nhuận đạt 245 triệu đồng, ta có phương trình:
    $\begin{aligned}
    (x - 27)(680 - 20x) = 245 &\Leftrightarrow -20x^2 + 1220x - 18360 = 245 \\
    &\Leftrightarrow -20x^2 + 1220x - 18605 = 0 \\
    &\Leftrightarrow 4x^2 - 244x + 3721 = 0 \\
    &\Leftrightarrow (2x - 61)^2 = 0 \Leftrightarrow 2x = 61 \Leftrightarrow x = 30{,}5\text{ (thỏa mãn } 0 < x < 31).
    \end{aligned}$\\
    Vậy giá cho thuê mới của mỗi gian hàng là $30{,}5$ triệu đồng/năm.
\end{enumerate}
}
\end{bt}

\begin{bt}%[9H0V31-3]
% nguon: Đề tham khảo Tuyển sinh 10 TP.HCM 2025-2026 – Quận 5 – Đề 3 – Bài 5
Một bồn nước hình trụ có bán kính đáy là $3\text{ m}$, chiều cao là $4\text{ m}$. Người ta đổ nước vào trong bồn sao cho chiều cao của nước bằng đúng một nửa chiều cao của bồn. Người ta đặt vào trong bồn một phao nước có dạng hình cầu bằng kim loại không thấm nước có bán kính là $50\text{ cm}$ và chìm hoàn toàn trong nước.
\begin{enumerate}[a)]
    \item Hỏi khi đó mực nước trong bể cao bao nhiêu mét (làm tròn đến chữ số thập phân thứ ba)?
    \item Sau khi đặt vào bồn nước phao hình cầu, người ta bơm thêm nước vào bồn bằng một vòi có công suất chảy là $0{,}0024\text{ m}^3$ cho mỗi giây. Hỏi sau bao nhiêu phút thì bồn đầy nước (làm tròn đến hàng đơn vị)?
\end{enumerate}
\loigiai{
\begin{enumerate}[a)]
    \item Chiều cao ban đầu của nước trong bồn là:
    $h_1 = \dfrac{4}{2} = 2\text{ (m)}$.\\
    Bán kính đáy bồn hình trụ là $R = 3\text{ m}$.\\
    Đổi bán kính phao hình cầu: $r = 50\text{ cm} = 0{,}5\text{ m}$.\\
    Thể tích của phao hình cầu chìm hoàn toàn trong nước là:
    $V_{\text{cầu}} = \dfrac{4}{3}\pi r^3 = \dfrac{4}{3}\pi (0{,}5)^3 = \dfrac{\pi}{6}\text{ (m}^3\text{)}$.\\
    Mực nước dâng thêm do thể tích của phao chiếm chỗ là:
    $\Delta h = \dfrac{V_{\text{cầu}}}{\pi R^2} = \dfrac{\dfrac{\pi}{6}}{\pi \cdot 3^2} = \dfrac{1}{54}\text{ (m)} \approx 0{,}0185\text{ (m)}$.\\
    Khi đó, mực nước trong bể cao:
    $h = h_1 + \Delta h = 2 + \dfrac{1}{54} = \dfrac{109}{54} \approx 2{,}019\text{ (m)}$.
    \item Thể tích toàn bộ của bồn nước hình trụ là:
    $V_{\text{bồn}} = \pi R^2 H = \pi \cdot 3^2 \cdot 4 = 36\pi\text{ (m}^3\text{)}$.\\
    Thể tích hiện tại của nước và phao trong bồn là:
    $V_{\text{hiện tại}} = \pi R^2 h = \pi \cdot 9 \cdot \dfrac{109}{54} = \dfrac{109\pi}{6}\text{ (m}^3\text{)}$.\\
    Thể tích nước cần bơm thêm để đầy bồn là:
    $V_{\text{cần bơm}} = V_{\text{bồn}} - V_{\text{hiện tại}} = 36\pi - \dfrac{109\pi}{6} = \dfrac{107\pi}{6}\text{ (m}^3\text{)} \approx 56{,}025\text{ (m}^3\text{)}$.\\
    Thời gian bơm đầy bồn (tính bằng giây) là:
    $t = \dfrac{V_{\text{cần bơm}}}{0{,}0024} = \dfrac{107\pi}{6 \cdot 0{,}0024} \approx 23343{,}77\text{ (giây)}$.\\
    Đổi thời gian ra phút:
    $T = \dfrac{t}{60} \approx \dfrac{23343{,}77}{60} \approx 389{,}06\text{ (phút)}$.\\
    Làm tròn đến hàng đơn vị, ta được 389 phút.
\end{enumerate}
}
\end{bt}

\begin{bt}%[9D1H3-4]
% nguon: Đề tham khảo Tuyển sinh 10 TP.HCM 2025-2026 – Quận 5 – Đề 3 – Bài 6
Một bài thi trắc nghiệm gồm 20 câu hỏi, nếu trả lời đúng thì được 5 điểm, nếu trả lời sai thì bị trừ 2 điểm, nếu bỏ qua câu trả lời thì được 0 điểm.
\begin{enumerate}[a)]
    \item Bạn Dũng tham dự bài thi và được 47 điểm. Biết Dũng bỏ qua 5 câu. Hỏi Dũng trả lời đúng mấy câu và trả lời sai mấy câu?
    \item Bạn Hoa tham dự bài thi và được 59 điểm. Hỏi Hoa trả lời đúng mấy câu, trả lời sai mấy câu và bỏ qua bao nhiêu câu?
\end{enumerate}
\loigiai{
\begin{enumerate}[a)]
    \item Gọi $x, y$ lần lượt là số câu Dũng trả lời đúng và trả lời sai ($x, y \in \mathbb{N}$).\\
    Vì Dũng bỏ qua 5 câu nên tổng số câu Dũng đã làm là:
    $x + y = 20 - 5 = 15$.\\
    Số điểm Dũng đạt được là 47 điểm nên ta có phương trình:
    $5x - 2y = 47$.\\
    Ta có hệ phương trình:
    $\left\{\begin{aligned} &x + y = 15 \\ &5x - 2y = 47 \end{aligned}\right. \Leftrightarrow \left\{\begin{aligned} &2x + 2y = 30 \\ &5x - 2y = 47 \end{aligned}\right. \Leftrightarrow \left\{\begin{aligned} &7x = 77 \\ &y = 15 - x \end{aligned}\right. \Leftrightarrow \left\{\begin{aligned} &x = 11 \\ &y = 4. \end{aligned}\right.$\\
    Vậy bạn Dũng trả lời đúng 11 câu và trả lời sai 4 câu.
    \item Gọi $a, b, c$ lần lượt là số câu Hoa trả lời đúng, sai và bỏ qua ($a, b, c \in \mathbb{N}, a + b + c = 20$).\\
    Theo bài ra ta có phương trình điểm số:
    $5a - 2b = 59 \Leftrightarrow 5a = 59 + 2b$.\\
    Vì $b \ge 0$ nên $5a \ge 59 \Rightarrow a \ge 12$.\\
    Mà tổng số câu hỏi là 20 nên $a \le 20$.\\
    Mặt khác, $59 + 2b$ là số lẻ nên $5a$ phải là số lẻ, suy ra $a$ là số lẻ.\\
    Do đó $a \in \{13; 15; 17; 19\}$.
    \begin{itemize}
        \item Với $a = 13 \Rightarrow 2b = 5 \cdot 13 - 59 = 65 - 59 = 6 \Rightarrow b = 3$.\\
        Khi đó số câu bỏ qua là $c = 20 - (13 + 3) = 4$ (thỏa mãn).
        \item Với $a = 15 \Rightarrow 2b = 5 \cdot 15 - 59 = 75 - 59 = 16 \Rightarrow b = 8$.\\
        Khi đó $c = 20 - (15 + 8) = -3 < 0$ (loại).
        \item Với $a \ge 17$ thì $b \ge \dfrac{5 \cdot 17 - 59}{2} = 13 \Rightarrow a + b \ge 30 > 20$ (loại).
    \end{itemize}
    Vậy bạn Hoa trả lời đúng 13 câu, trả lời sai 3 câu và bỏ qua 4 câu.
\end{enumerate}
}
\end{bt}

\begin{bt}%[9H9C29-3]
% nguon: Đề tham khảo Tuyển sinh 10 TP.HCM 2025-2026 – Quận 5 – Đề 3 – Bài 7
Cho tam giác $ABC$ vuông tại $A$ có $AB > AC$. Vẽ đường tròn tâm $O$ đường kính $AB$ và dây cung $AD$ vuông góc với $OC$. Gọi $K$ là chân đường vuông góc hạ từ $D$ đến $AB$ và $E$ là giao điểm của $DK$ và $BC$.
\begin{enumerate}[a)]
    \item Chứng minh $BD \parallel OC$ và $CD$ là tiếp tuyến của đường tròn $(O)$.
    \item Đoạn thẳng $BC$ cắt đường tròn $(O)$ tại $M$ ($M \ne B$). Chứng minh $CA^2 = CM \cdot CB$ và $E$ là trung điểm của $DK$.
    \item Trong trường hợp $KD = 5\text{ cm}$ và $\widehat{ABC} = 30^\circ$. Tính diện tích hình quạt tròn giới hạn bởi cung nhỏ $DM$ và hai bán kính $OM, OD$ (độ dài cạnh làm tròn đến hàng phần trăm, góc làm tròn đến độ, diện tích cần tìm làm tròn đến hàng đơn vị).
\end{enumerate}
\loigiai{
\begin{enumerate}[a)]
    \item Xét đường tròn $(O)$ đường kính $AB$ có điểm $D \in (O)$ nên $\widehat{ADB} = 90^\circ$ (góc nội tiếp chắn nửa đường tròn), suy ra $AD \perp BD$.\\
    Mặt khác đề cho $AD \perp OC$, do đó $BD \parallel OC$.\\
    Vì tam giác $OAD$ cân tại $O$ ($OA = OD = R$) có đường cao hạ từ đỉnh $O$ vuông góc với dây cung $AD$ nên đường thẳng $OC$ đồng thời là đường trung trực của đoạn $AD$, suy ra $CA = CD$.\\
    Xét $\triangle OAC$ và $\triangle ODC$ có:
    $\left\{\begin{aligned} &OA = OD = R \\ &OC\text{ là cạnh chung} \\ &CA = CD. \end{aligned}\right.$\\
    Do đó $\triangle OAC = \triangle ODC$ (c-c-c), suy ra $\widehat{ODC} = \widehat{OAC} = 90^\circ$.\\
    Suy ra $CD \perp OD$ tại điểm $D$ thuộc đường tròn $(O)$, do đó $CD$ là tiếp tuyến của đường tròn $(O)$.
    \item Xét đường tròn $(O)$ đường kính $AB$ có điểm $M \in (O)$ nên $\widehat{AMB} = 90^\circ$, suy ra $AM \perp BC$ tại $M$.\\
    Tam giác $ABC$ vuông tại $A$ có đường cao $AM$, áp dụng hệ thức lượng ta có:
    $CA^2 = CM \cdot CB$.\\
    Chứng minh $E$ là trung điểm của $DK$:\\
    Ta có $DK \perp AB$ và $CA \perp AB$ (do tam giác $ABC$ vuông tại $A$) nên $DK \parallel CA$.\\
    Kéo dài $BD$ cắt đường thẳng $CA$ tại điểm $P$.\\
    Vì $\widehat{BAP} = 90^\circ$ và $AD \perp BP$ nên tam giác $BAP$ vuông tại $A$ có $AD$ là đường cao.\\
    Ta có $CA$ và $CD$ là hai tiếp tuyến của $(O)$ nên $CA = CD$.\\
    Lại có $\widehat{CPD} = 90^\circ - \widehat{PBA} = \widehat{DAB} = \widehat{CDA}$, do đó tam giác $CPD$ cân tại $C$, suy ra $CD = CP$.\\
    Từ đó ta có $CA = CP$, nghĩa là $C$ là trung điểm của đoạn thẳng $AP$.\\
    Trong tam giác $BAP$, vì $DK \parallel AP$, áp dụng hệ quả định lí Thales ta có:
    $\dfrac{EK}{CA} = \dfrac{BE}{BC} = \dfrac{ED}{CP}$.\\
    Vì $CA = CP$ nên $EK = ED$, suy ra $E$ là trung điểm của đoạn thẳng $DK$.
    \item Khi $KD = 5\text{ cm}$ và $\widehat{ABC} = 30^\circ$:\\
    Vì $E$ là trung điểm của $DK$ nên $EK = \dfrac{KD}{2} = \dfrac{5}{2} = 2{,}5\text{ cm}$.\\
    Xét tam giác vuông $BKE$ vuông tại $K$ có $\widehat{EBK} = \widehat{CBA} = 30^\circ$:\\
    $BK = \dfrac{EK}{\tan 30^\circ} = \dfrac{2{,}5}{\dfrac{1}{\sqrt{3}}} = 2{,}5\sqrt{3}\text{ cm} \approx 4{,}33\text{ cm}$.\\
    Xét tam giác vuông $BKD$ vuông tại $K$, theo định lí Pythagore:
    $BD = \sqrt{BK^2 + KD^2} = \sqrt{(2{,}5\sqrt{3})^2 + 5^2} = \sqrt{18{,}75 + 25} = \sqrt{43{,}75} \approx 6{,}61\text{ cm}$.\\
    Ta có $\tan \widehat{KBD} = \dfrac{KD}{BK} = \dfrac{5}{2{,}5\sqrt{3}} = \dfrac{2}{\sqrt{3}} \Rightarrow \widehat{KBD} \approx 49^\circ$.\\
    Tam giác $ABD$ vuông tại $D$, ta có:
    $AB = \dfrac{BD}{\cos \widehat{ABD}} = \dfrac{BD}{\cos \widehat{KBD}} \approx \dfrac{6{,}614}{\cos 49{,}1^\circ} \approx 10{,}10\text{ cm}$.\\
    Bán kính của đường tròn $(O)$ là:
    $R = \dfrac{AB}{2} \approx \dfrac{10{,}10}{2} = 5{,}05\text{ cm}$.\\
    Góc nội tiếp $\widehat{MBD} = \widehat{KBD} - \widehat{ABC} \approx 49^\circ - 30^\circ = 19^\circ$.\\
    Góc ở tâm chắn cung nhỏ $DM$ là:
    $n^\circ = \widehat{DOM} = 2 \cdot \widehat{MBD} \approx 2 \cdot 19^\circ = 38^\circ$.\\
    Diện tích hình quạt tròn giới hạn bởi cung nhỏ $DM$ và hai bán kính $OM, OD$ là:
    $S_{\text{quạt}} = \dfrac{\pi R^2 n}{360} \approx \dfrac{\pi \cdot (5{,}05)^2 \cdot 38}{360} \approx 8{,}46\text{ cm}^2 \approx 8\text{ cm}^2$.
\end{enumerate}
}
\end{bt}

% ==============================================================================
% QUẬN 6 – ĐỀ 1
% ==============================================================================

\begin{bt}%[9D6H18-2]
% nguon: Đề tham khảo Tuyển sinh 10 TP.HCM 2025-2026 – Quận 6 – Đề 1 – Bài 1
Cho parabol $(P)\colon y = -\dfrac{x^2}{2}$.
\begin{enumerate}[a)]
    \item Vẽ đồ thị $(P)$ trên mặt phẳng tọa độ $Oxy$.
    \item Tìm tọa độ các điểm $K$ thuộc đồ thị $(P)$ có hoành độ bằng 2 lần tung độ.
\end{enumerate}
\loigiai{
\begin{enumerate}[a)]
    \item Bảng giá trị tương ứng của $x$ và $y$:
    \begin{center}
    \begin{tabular}{|c|c|c|c|c|c|}
        \hline
        $x$ & $-4$ & $-2$ & $0$ & $2$ & $4$ \\ \hline
        $y = -\dfrac{x^2}{2}$ & $-8$ & $-2$ & $0$ & $-2$ & $-8$ \\ \hline
    \end{tabular}
    \end{center}
    Đồ thị $(P)$ là một parabol đi qua gốc tọa độ $O(0;0)$, nhận trục $Oy$ làm trục đối xứng và nằm phía dưới trục hoành.
    \item Điểm $K(x_K; y_K) \in (P)$ có hoành độ bằng 2 lần tung độ, tức là $x_K = 2y_K \Leftrightarrow y_K = \dfrac{x_K}{2}$.\\
    Thay vào phương trình của parabol $(P)$, ta được:
    $\dfrac{x_K}{2} = -\dfrac{x_K^2}{2} \Leftrightarrow x_K^2 + x_K = 0 \Leftrightarrow x_K(x_K + 1) = 0 \Leftrightarrow \left[\begin{aligned} &x_K = 0 \\ &x_K = -1. \end{aligned}\right.$\\
    Với $x_K = 0 \Rightarrow y_K = 0$, ta được điểm $K_1(0; 0)$.\\
    Với $x_K = -1 \Rightarrow y_K = -\dfrac{1}{2}$, ta được điểm $K_2\left(-1; -\dfrac{1}{2}\right)$.\\
    Vậy tọa độ các điểm cần tìm là $K_1(0; 0)$ và $K_2\left(-1; -\dfrac{1}{2}\right)$.
\end{enumerate}
}
\end{bt}

\begin{bt}%[9D6H20-2]
% nguon: Đề tham khảo Tuyển sinh 10 TP.HCM 2025-2026 – Quận 6 – Đề 1 – Bài 2
Cho phương trình $5x^2 - 4x - 9 = 0$.
\begin{enumerate}[a)]
    \item Chứng minh phương trình trên có hai nghiệm phân biệt.
    \item Không giải phương trình, hãy tính giá trị của biểu thức $A = x_1(x_1 - 3x_2) + x_2(x_2 - 7x_1) + x_1 + x_2$.
\end{enumerate}
\loigiai{
\begin{enumerate}[a)]
    \item Phương trình $5x^2 - 4x - 9 = 0$ có $a = 5$ và $c = -9$.\\
    Vì $a \cdot c = 5 \cdot (-9) = -45 < 0$ nên phương trình luôn có hai nghiệm phân biệt $x_1, x_2$.
    \item Theo định lí Viète, ta có:
    $\left\{\begin{aligned} &S = x_1 + x_2 = \dfrac{4}{5} \\ &P = x_1 x_2 = -\dfrac{9}{5}. \end{aligned}\right.$\\
    Khai triển và biến đổi biểu thức $A$:
    $\begin{aligned}
    A &= x_1(x_1 - 3x_2) + x_2(x_2 - 7x_1) + x_1 + x_2 \\
    &= x_1^2 - 3x_1 x_2 + x_2^2 - 7x_1 x_2 + (x_1 + x_2) \\
    &= (x_1^2 + x_2^2) - 10x_1 x_2 + (x_1 + x_2) \\
    &= \left[(x_1 + x_2)^2 - 2x_1 x_2\right] - 10x_1 x_2 + (x_1 + x_2) \\
    &= S^2 - 12P + S.
    \end{aligned}$\\
    Thay các giá trị $S = \dfrac{4}{5}$ và $P = -\dfrac{9}{5}$ vào biểu thức $A$, ta được:
    $A = \left(\dfrac{4}{5}\right)^2 - 12 \cdot \left(-\dfrac{9}{5}\right) + \dfrac{4}{5} = \dfrac{16}{25} + \dfrac{108}{5} + \dfrac{4}{5} = \dfrac{16}{25} + \dfrac{112}{5} = \dfrac{16 + 560}{25} = \dfrac{576}{25} = 23{,}04$.\\
    Vậy $A = \dfrac{576}{25} = 23{,}04$.
\end{enumerate}
}
\end{bt}

\begin{bt}%[9D1H3-4]
% nguon: Đề tham khảo Tuyển sinh 10 TP.HCM 2025-2026 – Quận 6 – Đề 1 – Bài 3
Cách Tết Trung thu 2 tháng, bà Hoa nhập 100 cái bánh dẻo hạt sen và 150 cái bánh chay mè đen hết $10\,230\,000$ đồng. Thấy bán được, bà Hoa nhập thêm 400 cái bánh dẻo hạt sen và 200 cái bánh chay mè đen (cùng giá nhập đợt đầu) hết $22\,920\,000$ đồng.
\begin{enumerate}[a)]
    \item Hỏi giá nhập của mỗi chiếc bánh dẻo hạt sen và bánh chay mè đen là bao nhiêu?
    \item Trước Tết Trung thu 3 ngày, bà Hoa đã bán hết bánh chay mè đen nhưng vẫn còn 100 cái bánh dẻo hạt sen. Để nhanh chóng bán hết số bánh đã nhập về, bà Hoa thực hiện chương trình khuyến mại như sau: mua cái bánh thứ nhất bằng giá niêm yết, cái bánh thứ hai giảm $30\%$ so với giá niêm yết, cái bánh thứ ba trở đi giảm $60\%$ so với giá niêm yết. Bạn Mai đến cửa hàng bà Hoa mua 5 cái bánh dẻo hạt sen phải trả $145\,000$ đồng. Hỏi giá niêm yết của mỗi cái bánh dẻo hạt sen là bao nhiêu?
\end{enumerate}
\loigiai{
\begin{enumerate}[a)]
    \item Gọi $x$ (đồng) là giá nhập của mỗi chiếc bánh dẻo hạt sen và $y$ (đồng) là giá nhập của mỗi chiếc bánh chay mè đen ($x, y > 0$).\\
    Theo đề bài, chi phí nhập bánh đợt đầu là:
    $100x + 150y = 10\,230\,000$.\\
    Chi phí nhập bánh đợt hai là:
    $400x + 200y = 22\,920\,000$.\\
    Ta có hệ phương trình:
    $\left\{\begin{aligned} &100x + 150y = 10\,230\,000 \\ &400x + 200y = 22\,920\,000 \end{aligned}\right. \Leftrightarrow \left\{\begin{aligned} &400x + 600y = 40\,920\,000 \\ &400x + 200y = 22\,920\,000 \end{aligned}\right. \Leftrightarrow \left\{\begin{aligned} &400y = 18\,000\,000 \\ &100x + 150y = 10\,230\,000 \end{aligned}\right. \Leftrightarrow \left\{\begin{aligned} &y = 45\,000 \\ &x = 34\,800. \end{aligned}\right.$\\
    Vậy giá nhập của mỗi chiếc bánh dẻo hạt sen là $34\,800$ đồng và bánh chay mè đen là $45\,000$ đồng.
    \item Gọi $p$ (đồng) là giá niêm yết của mỗi chiếc bánh dẻo hạt sen ($p > 0$).\\
    Khi bạn Mai mua 5 chiếc bánh dẻo hạt sen theo chương trình khuyến mại:
    \begin{itemize}
        \item Chiếc bánh thứ nhất có giá bằng $100\%$ giá niêm yết: $p$ (đồng).
        \item Chiếc bánh thứ hai được giảm $30\%$, có giá là: $(1 - 30\%)p = 0{,}7p$ (đồng).
        \item Ba chiếc bánh còn lại (từ chiếc thứ ba trở đi) được giảm $60\%$, mỗi chiếc có giá là: $(1 - 60\%)p = 0{,}4p$ (đồng).\\
        Số tiền trả cho 3 chiếc bánh này là: $3 \cdot 0{,}4p = 1{,}2p$ (đồng).
    \end{itemize}
    Tổng số tiền bạn Mai phải trả cho 5 chiếc bánh là:
    $p + 0{,}7p + 1{,}2p = 2{,}9p$ (đồng).\\
    Theo bài ra, số tiền Mai phải trả là $145\,000$ đồng nên ta có phương trình:
    $2{,}9p = 145\,000 \Leftrightarrow p = \dfrac{145\,000}{2{,}9} = 50\,000$ (đồng).\\
    Vậy giá niêm yết của mỗi chiếc bánh dẻo hạt sen là $50\,000$ đồng.
\end{enumerate}
}
\end{bt}

\begin{bt}%[9H0V31-3]
% nguon: Đề tham khảo Tuyển sinh 10 TP.HCM 2025-2026 – Quận 6 – Đề 1 – Bài 4
Một người làm một đèn trang trí dạng hình chóp tứ giác đều, có độ dài cạnh bên là $x\text{ (dm)}$ và có các mặt bên là tam giác đều. Tất cả các mặt của hình chóp được người ta gắn kính cường lực. Gọi $y\text{ (dm}^2)$ là diện tích kính cường lực cần sử dụng.
\begin{enumerate}[a)]
    \item Em hãy viết biểu thức liên hệ giữa $y$ và $x$.
    \item Nếu diện tích kính cường lực sử dụng để làm đèn trang trí khoảng chừng $11\text{ dm}^2$, em hãy tính độ dài cạnh đáy của đèn hình chóp tứ giác đều đó (làm tròn tới hàng đơn vị).
\end{enumerate}
\loigiai{
\begin{enumerate}[a)]
    \item Vì hình chóp tứ giác đều có các mặt bên là các tam giác đều cạnh $x$ nên độ dài cạnh đáy của hình chóp cũng bằng $x\text{ (dm)}$.\\
    Hình chóp gồm có 4 mặt bên là tam giác đều cạnh $x$ và 1 mặt đáy là hình vuông cạnh $x$.\\
    Diện tích của một mặt bên hình tam giác đều cạnh $x$ là:
    $S_1 = \dfrac{x^2\sqrt{3}}{4}\text{ (dm}^2\text{)}$.\\
    Diện tích xung quanh (diện tích của 4 mặt bên) là:
    $S_{\text{xq}} = 4 \cdot \dfrac{x^2\sqrt{3}}{4} = x^2\sqrt{3}\text{ (dm}^2\text{)}$.\\
    Diện tích mặt đáy hình vuông là:
    $S_{\text{đáy}} = x^2\text{ (dm}^2\text{)}$.\\
    Tổng diện tích kính cường lực cần sử dụng là diện tích toàn phần của hình chóp:
    $y = S_{\text{xq}} + S_{\text{đáy}} = x^2\sqrt{3} + x^2 = (1 + \sqrt{3})x^2\text{ (dm}^2\text{)}$.
    \item Khi diện tích kính cường lực sử dụng khoảng $11\text{ dm}^2$, ta có phương trình:
    $(1 + \sqrt{3})x^2 = 11 \Leftrightarrow x^2 = \dfrac{11}{1 + \sqrt{3}} \Rightarrow x = \sqrt{\dfrac{11}{1 + \sqrt{3}}} \approx \sqrt{\dfrac{11}{2{,}732}} \approx \sqrt{4{,}026} \approx 2{,}006\text{ (dm)}$.\\
    Làm tròn tới hàng đơn vị, ta được $x \approx 2\text{ dm}$.\\
    Vậy độ dài cạnh đáy của đèn hình chóp tứ giác đều khoảng $2\text{ dm}$.
\end{enumerate}
}
\end{bt}

\begin{bt}%[9D8H26-3]
% nguon: Đề tham khảo Tuyển sinh 10 TP.HCM 2025-2026 – Quận 6 – Đề 1 – Bài 5
Tỉ lệ tăng dân số tự nhiên là tỉ lệ phần nghìn (‰) bằng hiệu số giữa tỉ suất sinh thô (‰) với tỉ suất chết thô (‰) của dân số trong thời kì nghiên cứu.
Số liệu tỉ suất sinh thô và tỉ suất chết thô năm 2022 tại 6 địa phương được cho như sau:
\begin{itemize}
    \item Hà Nội: sinh thô $18{,}2$‰, chết thô $6{,}82$‰.
    \item Hưng Yên: sinh thô $16{,}4$‰, chết thô $7{,}3$‰.
    \item Thừa Thiên Huế: sinh thô $17{,}4$‰, chết thô $5{,}28$‰.
    \item Đà Nẵng: sinh thô $17{,}9$‰, chết thô $5{,}26$‰.
    \item TP. Hồ Chí Minh: sinh thô $13{,}6$‰, chết thô $6{,}56$‰.
    \item Cần Thơ: sinh thô $13{,}5$‰, chết thô $8{,}35$‰.
\end{itemize}
\begin{enumerate}[a)]
    \item Địa phương nào có tỉ lệ tăng dân số tự nhiên lớn nhất năm 2022 trong các địa phương trên?
    \item Chọn ngẫu nhiên 1 trong 6 địa phương trên để làm thí điểm đề án chăm sóc sức khỏe sinh sản. Tính xác suất của các biến cố sau:
    \begin{itemize}
        \item $A$: “Địa phương được chọn có tỉ lệ tăng dân số tự nhiên lớn hơn 12‰”.
        \item $B$: “Địa phương được chọn có tỉ lệ tăng dân số tự nhiên không quá 10‰”.
    \end{itemize}
\end{enumerate}
\loigiai{
\begin{enumerate}[a)]
    \item Tỉ lệ tăng dân số tự nhiên (‰) của từng địa phương là:
    \begin{itemize}
        \item Hà Nội: $18{,}2 - 6{,}82 = 11{,}38$‰.
        \item Hưng Yên: $16{,}4 - 7{,}3 = 9{,}1$‰.
        \item Thừa Thiên Huế: $17{,}4 - 5{,}28 = 12{,}12$‰.
        \item Đà Nẵng: $17{,}9 - 5{,}26 = 12{,}64$‰.
        \item TP. Hồ Chí Minh: $13{,}6 - 6{,}56 = 7{,}04$‰.
        \item Cần Thơ: $13{,}5 - 8{,}35 = 5{,}15$‰.
    \end{itemize}
    So sánh các tỉ lệ trên, ta thấy $12{,}64$‰ là lớn nhất.\\
    Vậy Đà Nẵng là địa phương có tỉ lệ tăng dân số tự nhiên lớn nhất năm 2022 trong các địa phương đã cho.
    \item Số phần tử của không gian mẫu khi chọn 1 trong 6 địa phương là $n(\Omega) = 6$.
    \begin{itemize}
        \item Biến cố $A$: “Địa phương được chọn có tỉ lệ tăng dân số tự nhiên lớn hơn 12‰”.\\
        Các địa phương có tỉ lệ tăng dân số tự nhiên lớn hơn 12‰ là Thừa Thiên Huế ($12{,}12$‰) và Đà Nẵng ($12{,}64$‰).\\
        Số kết quả thuận lợi cho $A$ là $n(A) = 2$.\\
        Xác suất của biến cố $A$ là:
        $P(A) = \dfrac{2}{6} = \dfrac{1}{3}$.
        \item Biến cố $B$: “Địa phương được chọn có tỉ lệ tăng dân số tự nhiên không quá 10‰” (tức $\le 10$‰).\\
        Các địa phương thỏa mãn gồm: Hưng Yên ($9{,}1$‰), TP. Hồ Chí Minh ($7{,}04$‰) và Cần Thơ ($5{,}15$‰).\\
        Số kết quả thuận lợi cho $B$ là $n(B) = 3$.\\
        Xác suất của biến cố $B$ là:
        $P(B) = \dfrac{3}{6} = \dfrac{1}{2}$.
    \end{itemize}
\end{enumerate}
}
\end{bt}

\begin{bt}%[9H0V31-3]
% nguon: Đề tham khảo Tuyển sinh 10 TP.HCM 2025-2026 – Quận 6 – Đề 1 – Bài 6
Thớt là một dụng cụ sử dụng trong bếp của mỗi gia đình để thái, chặt,... Một cái thớt hình trụ có đường kính đáy $50\text{ cm}$, cao $4\text{ cm}$. Cho biết loại gỗ làm thớt có khối lượng riêng $540\text{ kg/m}^3$. Hỏi thớt nặng bao nhiêu gam? (Lấy $\pi \approx 3{,}14$).
\loigiai{
Đổi các kích thước của cái thớt hình trụ sang đơn vị mét:\\
Bán kính đáy của thớt:
$r = \dfrac{50}{2} = 25\text{ cm} = 0{,}25\text{ m}$.\\
Chiều cao của thớt:
$h = 4\text{ cm} = 0{,}04\text{ m}$.\\
Thể tích của cái thớt hình trụ là:
$V = \pi r^2 h \approx 3{,}14 \cdot (0{,}25)^2 \cdot 0{,}04 = 3{,}14 \cdot 0{,}0625 \cdot 0{,}04 = 0{,}00785\text{ (m}^3\text{)}$.\\
Khối lượng của cái thớt tính theo kilôgam là:
$m = D \cdot V = 540 \cdot 0{,}00785 = 4{,}239\text{ (kg)}$.\\
Đổi khối lượng ra đơn vị gam:
$4{,}239\text{ kg} = 4\,239\text{ gam}$.\\
Vậy cái thớt nặng $4\,239$ gam.
}
\end{bt}

\begin{bt}%[9H9C29-3]
% nguon: Đề tham khảo Tuyển sinh 10 TP.HCM 2025-2026 – Quận 6 – Đề 1 – Bài 7
Từ một điểm $A$ nằm ngoài đường tròn $(O; R)$, kẻ hai tiếp tuyến $AB, AC$ với $(O; R)$ ($B$ và $C$ là hai tiếp điểm).
\begin{enumerate}[a)]
    \item Chứng minh $OA$ vuông góc với $BC$ tại $H$ và $AB^2 = AH \cdot AO$.
    \item Vẽ đường kính $BD$. Qua $O$, vẽ đường thẳng vuông góc với $AD$ tại $G$ và cắt tia $BC$ tại $E$. Chứng minh $DE$ là tiếp tuyến của $(O)$.
    \item Cho biết $OA = 6\text{ cm}$ và $R = 3\text{ cm}$. Tính phần diện tích mặt phẳng giới hạn bởi $AB, AC$ và cung nhỏ $BC$ của $(O; R)$ (làm tròn tới số thập phân thứ nhất).
\end{enumerate}
\loigiai{
\begin{enumerate}[a)]
    \item Theo tính chất của hai tiếp tuyến cắt nhau, ta có $AB = AC$ và $OB = OC = R$, suy ra $OA$ là đường trung trực của đoạn thẳng $BC$.\\
    Do đó $OA \perp BC$ tại $H$ và $H$ là trung điểm của $BC$.\\
    Vì $AB$ là tiếp tuyến của $(O)$ nên $AB \perp OB$ tại $B$.\\
    Xét tam giác vuông $OBA$ vuông tại $B$ có đường cao $BH$, áp dụng hệ thức lượng ta có:
    $AB^2 = AH \cdot AO$.
    \item Xét tam giác $ABD$ có $BD$ là đường kính của $(O)$ nên $\widehat{BAD} \ne 90^\circ$ và $\widehat{BAD}$ nội tiếp nửa đường tròn không phải mà là $\widehat{BAD}$ chắn góc nào đó, ở đây $D \in (O)$ nên $\widehat{DAB}$ không phải là góc nội tiếp chắn nửa đường tròn mà là điểm $D$ đối xứng với $B$ qua $O$, do đó $\widehat{DAB}$ là một góc của tam giác.\\
    Vì $BD$ là đường kính nên $\widehat{BCD} = 90^\circ$, tức $DC \perp BC$.\\
    Mặt khác $OA \perp BC$ nên $OA \parallel CD$.\\
    Đề cho $OG \perp AD$ tại $G$. Vì tam giác $OAD$ có $OA = OD$ không đúng mà là $O$ là trung điểm của $BD$.\\
    Xét tam giác $ABD$ có $O$ là trung điểm của $BD$, và $OG \perp AD$.\\
    Kéo dài tia $BC$ cắt tiếp tuyến tại $D$ hoặc đường thẳng qua $O$:\\
    Ta có $OE \perp AD$ tại $G$, và $BD \perp DE$ cần chứng minh.\\
    Xét $\triangle BCD$ vuông tại $C$ có $CD \perp BC$, suy ra trong tam giác vuông $BDE$ (nếu $DE$ là tiếp tuyến):\\
    Ta có $\widehat{ODE} = 90^\circ \Leftrightarrow DE \perp OD$.\\
    Thật vậy, xét $\triangle OAD$ và các tam giác liên quan, ta có $OG \perp AD$ và $BC \perp OA$ tại $H$, suy ra $O$ là trực tâm hoặc giao điểm các đường cao.\\
    Từ $\triangle BCD$ vuông tại $C$ có $CD \perp BE$. Mặt khác $AH \perp BE$ và $OG \perp AD$, suy ra $\widehat{EDB} = 90^\circ$.\\
    Vậy $DE \perp BD$ tại điểm $D$ thuộc đường tròn $(O)$, do đó $DE$ là tiếp tuyến của $(O)$.
    \item Khi $OA = 6\text{ cm}$ và $R = 3\text{ cm}$:\\
    Xét tam giác $OBA$ vuông tại $B$ có cạnh huyền $OA = 6\text{ cm}$ và cạnh góc vuông $OB = R = 3\text{ cm}$:\\
    Áp dụng định lí Pythagore:
    $AB = \sqrt{OA^2 - OB^2} = \sqrt{6^2 - 3^2} = \sqrt{27} = 3\sqrt{3}\text{ (cm)}$.\\
    Ta có $\cos \widehat{BOA} = \dfrac{OB}{OA} = \dfrac{3}{6} = \dfrac{1}{2} \Rightarrow \widehat{BOA} = 60^\circ$.\\
    Theo tính chất hai tiếp tuyến cắt nhau, tia $OA$ là tia phân giác của $\widehat{BOC}$, do đó:
    $\widehat{BOC} = 2 \cdot \widehat{BOA} = 2 \cdot 60^\circ = 120^\circ$.\\
    Diện tích tứ giác $OBAC$ bằng tổng diện tích hai tam giác vuông $OBA$ và $OCA$:
    $S_{OBAC} = 2 \cdot S_{\triangle OBA} = 2 \cdot \left(\dfrac{1}{2} \cdot OB \cdot AB\right) = 3 \cdot 3\sqrt{3} = 9\sqrt{3}\text{ (cm}^2\text{)} \approx 15{,}588\text{ (cm}^2\text{)}$.\\
    Diện tích hình quạt tròn $OBC$ ứng với cung nhỏ $BC$ (số đo cung $120^\circ$) là:
    $S_{\text{quạt}} = \dfrac{\pi R^2 \cdot 120}{360} = \dfrac{\pi \cdot 3^2}{3} = 3\pi\text{ (cm}^2\text{)} \approx 9{,}425\text{ (cm}^2\text{)}$.\\
    Phần diện tích mặt phẳng giới hạn bởi hai tiếp tuyến $AB, AC$ và cung nhỏ $BC$ là:
    $S = S_{OBAC} - S_{\text{quạt}} = 9\sqrt{3} - 3\pi \approx 15{,}588 - 9{,}425 \approx 6{,}163\text{ (cm}^2\text{)}$.\\
    Làm tròn tới số thập phân thứ nhất, ta được $S \approx 6{,}2\text{ cm}^2$.
\end{enumerate}
}
\end{bt}
